\documentclass[11pt,reqno]{amsart}
\usepackage[margin=1.15in]{geometry}
\usepackage{amsmath,amssymb,amsthm,mathtools}
\usepackage{booktabs}
\usepackage[colorlinks=true,linkcolor=blue,citecolor=blue,urlcolor=blue]{hyperref}

\theoremstyle{plain}
\newtheorem{theorem}{Theorem}[section]
\newtheorem{proposition}[theorem]{Proposition}
\newtheorem{lemma}[theorem]{Lemma}
\newtheorem{corollary}[theorem]{Corollary}
\newtheorem{conjecture}[theorem]{Conjecture}
\theoremstyle{definition}
\newtheorem{computation}[theorem]{Computation}
\newtheorem{observation}[theorem]{Observation}
\newtheorem{remark}[theorem]{Remark}

\DeclareMathOperator{\Diff}{Diff}
\DeclareMathOperator{\ord}{ord}
\DeclareMathOperator{\disc}{disc}
\DeclareMathOperator{\tr}{tr}
\DeclareMathOperator{\nrd}{nrd}
\DeclareMathOperator{\trd}{trd}
\DeclareMathOperator{\Jac}{Jac}
\newcommand{\Q}{\mathbb{Q}}
\newcommand{\Z}{\mathbb{Z}}
\newcommand{\C}{\mathbb{C}}
\newcommand{\R}{\mathbb{R}}
\newcommand{\F}{\mathbb{F}}
\newcommand{\HH}{\mathbb{H}}
\newcommand{\Ldual}{L^{\vee}}
\newcommand{\Lm}{L_{-}}
\newcommand{\Lp}{L_{+}}

\begin{document}

\title[Evaluating forms on Shimura curves by the Shimizu lift]
{Evaluating weight-two forms on Shimura curves by the Shimizu lift,\\
and what the values decide}

\author{}
\date{}

\begin{abstract}
We record how the differentials of an Atkin--Lehner quotient $X_0^D(N)/W$ are evaluated at
points of the upper half-plane by Nelson's formula for the Shimizu lift, how the formula is
extended from the newforms of level exactly $DN$ to the whole space of differentials, and the
uses made of the values: Noether's count of quadrics through the canonical image, which decides
hyperellipticity numerically; in genus three with an Atkin--Lehner involution whose quotient has
genus one, an explicit hyperelliptic model recognised from three rational invariants and confirmed
by point counts; and in genus two an equation over $\Q$ obtained from the differentials alone,
where the sextic is a linear system in its coefficients, the coordinate is fixed by rational CM
points, and the counts are needed only for the quadratic twist. The derivative of a form that the
last two require is taken inside the lattice sum rather than by a contour, which shares the
enumeration of lattice points between a value and its derivative. We state precisely what each
step proves and what it does not, and we record the three traps met on the way: the working
precision must rise with the cutoff of the lattice sum, the Eichler order that Magma returns is not
the same from one process to the next, and a rational reconstruction must take its denominator
bound from the accuracy achieved rather than from the working precision.
\end{abstract}

\maketitle

\section{The formula}

Let $B$ be the indefinite quaternion algebra over $\Q$ of discriminant $D$, let $R$ be an Eichler
order of level $N$ with $(D,N)=1$ and $DN$ squarefree, and let $X=X_0^D(N)$ be the Shimura curve
$\Gamma\backslash\HH$ for the group $\Gamma$ of units of reduced norm one, embedded in
$\mathrm{SL}_2(\R)$ by a fixed splitting $\iota$. A weight-$k$ form $f_B$ on $X$ with
Jacquet--Langlands correspondent a newform $f=\sum a_f(n)q^n$ of level exactly $DN$, $a_f(1)=1$,
is given by the following theorem of Nelson \cite[Theorem~3.1]{Nelson}, which is the Shimizu
lift made explicit. For $z\in\HH$ put $\sigma_z=\begin{psmallmatrix} y^{1/2} & x y^{-1/2}\\ 0 &
y^{-1/2}\end{psmallmatrix}$, so that $\sigma_z i = z$, and for a real $2\times 2$ matrix put
\[
X\!\begin{pmatrix} a & b\\ c & d\end{pmatrix} = \frac{a+d+i(b-c)}{2},
\qquad
P\!\begin{pmatrix} a & b\\ c & d\end{pmatrix} = \frac{a^2+b^2+c^2+d^2}{2}.
\]
For $d\mid DN$ let $R^{(d)}=R+\tfrac{DN}{d}R^\vee$, where $R^\vee$ is the dual of $R$ for the
trace form; it is a lattice between $R$ and $R^\vee$ of index $d^2$ over $R$, and the reduced
norms of its elements have denominators dividing $d$. Then for $z_1,z_2\in\HH$,
\begin{equation}\label{eq:nelson}
(y_1y_2)^{k/2}\,\overline{f_B(z_1)}\,f_B(z_2)
= \sum_{d\mid DN}\frac{\mu((d,N))}{d\,a_f(d)}
\sum_{\substack{\alpha\in R^{(d)}\\ \nrd\alpha>0}}
a_f\big(d\,\nrd\alpha\big)\,e^{ik\arg X(\alpha_{12})}\,W_k\big(4\pi\lvert X(\alpha_{12})\rvert\big),
\end{equation}
where $\alpha_{12}=\sigma_{z_1}^{-1}\iota(\alpha)\sigma_{z_2}$ and
\[
W_k(x)=2^{1-k}\sum_{n\ge 1} n^k\big(nx\,K_{k-1}(nx)-K_k(nx)\big).
\]
Two features of \eqref{eq:nelson} govern everything below.

First, the left side is a product of a value and a conjugate value. An individual value
$f_B(z)$ is defined only up to a constant of modulus one, and the right side determines the
product. Fixing $z_1$ once and reading $\overline{f_B(z_1)}f_B(z_j)$ across $j$ gives every form
with its own constant in front, which for the purposes below (a basis of differentials up to
rescaling of its members) costs nothing.

Second, the sum is finite once truncated, and the truncation is an enumeration of a positive
definite quaternary form. Since $2\lvert X(\alpha)\rvert^2=P(\alpha)+\det\alpha$, the condition
$4\pi\lvert X\rvert\le c$ is implied by $P(\alpha)\le 2(c/4\pi)^2$, and $P$ is positive definite
on $M_2(\R)$. The function $W_k$ decays like $e^{-x}$, so the terms dropped at cutoff $c$ are of
size about $e^{-c}$ relative to the largest; in practice the cutoff buys about $0.4$ decimal
digits per unit, with a scatter of a digit or so between bases and between points, and no trend
in $DN$. There is no fundamental domain, no reduction and no sampling anywhere in the method.

The ingredients are therefore: the algebra and an Eichler order, a real splitting, the lattices
$R^{(d)}$, the Fourier coefficients $a_f(n)$ for $n\le DN(c/4\pi)^2$, and $W_k$. The only
transcendental ingredient is $W_k$, evaluated once per lattice point; it is tabulated once per
cutoff by a Chebyshev approximation of $W_k(x)e^x$ on panels in geometric progression, the
geometric spacing being what keeps the singularity of $W_k$ at the origin equally far from every
panel.

\subsection*{The hypothesis of the theorem}
Theorem~3.1 is stated for a newform of level exactly $DN$. The space of differentials of
$X_0^D(N)$ is the $D$-new part of $S_2(\Gamma_0(DN))$, which also contains, for every proper
divisor $M$ of $N$, the oldclass of each newform of level $DM$. Whether the newforms of level
exactly $DN$ span the $W$-fixed differentials of a given quotient is a dimension count that takes
a second, and it fails at a good many quotients of interest. The next section removes the
restriction.

\section{The oldform extension}

Fix $M\mid N$, let $T$ be the set of primes dividing $N/M$, and let $g$ be a newform of level
exactly $DM$ with correspondent $g_B$ on $X_0^D(M)$. Let $R_N\subset R_M$ be Eichler orders of
levels $N$ and $M$, $R_M=R_N+M R_1$ for a maximal order $R_1\supset R_N$. For $p\in T$ let $w_p$
be an element of $R_N$ of reduced norm $p$ normalising $R_N$ (an Atkin--Lehner element), and for
$S\subseteq T$ let $w_S=\prod_{p\in S}w_p$. The weight-two slash operator
\[
(f\mid w)(z)=\nrd(w)\,j(w,z)^{-2}f(wz)
\]
is the pullback of the differential $f(z)\,dz$ under $z\mapsto wz$, and the oldclass of $g$ at
level $N$ is spanned by the $2^{\#T}$ pullbacks $g_B\mid w_S$: $S=\emptyset$ is the natural
projection $X_0^D(N)\to X_0^D(M)$ and $S=\{p\}$ the second degeneracy map at $p$. Any two choices
of $w_p$ differ by a unit of $R_N$ of reduced norm one, up to sign, so $g_B\mid w_S$ does not
depend on the choice.

\begin{proposition}\label{prop:oldclass}
On the oldclass of $g$ the Atkin--Lehner group acts as follows.
\begin{enumerate}
\item For $q\mid DM$, $w_q$ acts on every $g_B\mid w_S$ by the quaternionic eigenvalue
$\varepsilon_q(g)$ of $g_B$ on $X_0^D(M)$, which is $a_g(q)$ for $q\mid D$ and $-a_g(q)$ for
$q\mid M$.
\item For $p\in T$, $w_p$ sends $g_B\mid w_S$ to $g_B\mid w_{S\,\triangle\,\{p\}}$ with no scalar.
\end{enumerate}
Consequently the eigenvectors are $v_\chi=\sum_{S\subseteq T}\chi(S)\,g_B\mid w_S$ for the
characters $\chi$ of $\{\pm1\}^T$, and $w_m$ acts on $v_\chi$ by
$\prod_{q\mid m,\ q\mid DM}\varepsilon_q(g)\prod_{p\mid m,\ p\in T}\chi(p)$.
\end{proposition}

\begin{proof}
For (1), an Atkin--Lehner element $w_q$ of $R_N$ at a prime $q\mid DM$ lies in $R_M$ and
normalises it, the two orders agreeing locally at $q$, so it is also an Atkin--Lehner element at
level $M$ and acts on $g_B$ by its eigenvalue there; and $w_qw_S$ and $w_Sw_q$ generate the same
two-sided ideal of $R_N$, so they differ by a unit of $R_N$ of reduced norm one, on which the
slash operator is trivial. That eigenvalue is $\varepsilon_q(g)=a_g(q)$ for $q\mid D$ and
$-a_g(q)$ for $q\mid M$: the classical Atkin--Lehner eigenvalue of a newform at a prime exactly
dividing its level is $-a_g(q)$, and the Jacquet--Langlands correspondence reverses the sign at
the primes where the quaternion algebra ramifies \cite{JL,Voight}.
For (2), $w_p^2$ has reduced norm $p^2$ and lies in $R_N$, and $w_p^2/p$ is a unit of $R_N$ of
reduced norm one (this is checked exactly in the code for every $w_p$ used). The weight-two slash
is trivial on scalars, by the factor $\nrd(w)$, and on units of norm one, which act trivially on
forms for $\Gamma$.
\end{proof}

The dimension counts of the next paragraph, and the direct measurement of each column's
eigenvalues, are consistency checks on the convention rather than its justification.

The $W$-fixed differentials of $X_0^D(N)/W$ are then the $W$-fixed newforms of level exactly
$DN$ together with one $v_\chi$ for each pair $(g,\chi)$ whose eigenvalue product is $1$ for every
$m\in W$. This count must equal the genus, and the code refuses to continue when it does not;
agreement at every quotient tried, including classes with two primes in $T$, is the check of the
eigenvalue bookkeeping, and in particular of the sign at $q\mid M$.

\subsection*{Evaluating a pullback}
Formula \eqref{eq:nelson} at level $M$ gives $S_g(z_1,z_2)=y_1y_2\,\overline{g_B(z_1)}\,g_B(z_2)$.
Since $\operatorname{Im}(wz)=\nrd(w)\,y/\lvert j(w,z)\rvert^2$,
\begin{equation}\label{eq:pullback}
y_1y_2\,\overline{g_B(z_1)}\,(g_B\mid w_S)(z_2)=\frac{\overline{j(w_S,z_2)}}{j(w_S,z_2)}\,
S_g(z_1,w_Sz_2).
\end{equation}
So every member of the oldclass is the level-$M$ evaluator at the moved point, times a phase,
with the \emph{same} constant $\overline{g_B(z_1)}$ in front for every $S$. Nothing is left to
normalise: the relative scale of the pullbacks is fixed by the slash operator, and that is what
makes $v_\chi$ an eigenvector with eigenvalue exactly $\chi(p)$ rather than some $c$ and $1/c$.
The sign and scale are also measured directly: for each generator $m$ of $W$ and a few points
$z$, the ratio of the value of every column at $w_mz$, times $\overline{j}/j$, to its value at
$z$ is compared with the predicted eigenvalue, the rejected characters of each oldclass being
carried along so that their predicted eigenvalues $-1$ are measured too. On every control the
ratios are $\pm1$ to the working accuracy.

\subsection*{The Atkin--Lehner elements}
$w_m$ is found in the two-sided ideal $J_m=\{x\in R_N:\trd(xy)\in m\Z\text{ for all }y\in R_N\}$,
the radical of $R_N/mR_N$, of index $m^2$, as an element of reduced norm exactly $m$; such an
element generates $J_m$ and hence normalises $R_N$, which is checked directly as well. It is
produced as a short vector of $J_m$ for the positive definite majorant of the real splitting.
Magma's normaliser routines are for definite orders, and its principal-generator routine can
return a generator of negative norm, so neither is used.

\section{Sample points and cost}

The values are taken at points of $\HH$ spread by the square root of the arithmetic volume of
$\Gamma$. The spread matters for the count of the next section: rows from nearby points are
nearly proportional, and a matrix of quadratic monomials of effective rank near the number of
columns always has a quadric through it. The symptom of clustered points is a smooth ramp of
singular values instead of a gap, and a verdict is refused unless the gap reaches a stated size.

The cost is the lattice enumeration. The number of lattice points summed is, to within a few
percent at every base measured,
\[
10\,\Big(\frac{c}{4\pi}\Big)^4\,\frac{\sigma_2(DN)}{DN},\qquad
\frac{\sigma_2(DN)}{DN}=\prod_{p\mid DN}\Big(p+\frac1p\Big)=DN\prod_{p\mid DN}\Big(1+\frac1{p^2}\Big),
\]
so the cost of one evaluation is linear in $DN$ at a fixed cutoff, with a constant that is
larger when $DN$ has small prime factors, and grows like the fourth power of the cutoff. Several
eigenforms cost one enumeration, not several: everything but $a_f(n)$ is shared, so a basis of
$g$ differentials at a pair of points is one enumeration. The Fourier coefficients are needed up
to $DN(c/4\pi)^2$, and the newform decomposition at level $DN$ is the first ingredient to become
expensive, well beyond the levels of interest here.

\section{Noether's count}

Let $X$ have genus $g\ge3$ and let $u_1,\dots,u_g$ be a basis of its differentials. The canonical
image lies on $(g-2)(g-3)/2$ independent quadrics if $X$ is not hyperelliptic (Noether), and on
$(g-1)(g-2)/2$ if it is, the image being then a rational normal curve. At genus three this is
$0$ against $1$, at genus four $1$ against $3$. With values of the basis at $m$ points the
quadrics through the image are the kernel of the $m\times\binom{g+1}{2}$ matrix of quadratic
monomials, and the count is read off its singular values: rows and columns are normalised, the
matrix is realified, and the number of relations is the position of the largest ratio between
consecutive singular values, accepted only when that ratio is at least $10^4$ and the smallest
singular value is below $10^{-4}$; otherwise, when nothing is near zero, there is no relation, and
when the spectrum is a ramp the count is declared inconclusive.

\subsection*{What a numerical kernel dimension proves}
Nothing, strictly: a small singular value is a small number. What the reading relies on is a
separation that has to be established in two directions, and the second one is easy to omit.
A residual at the truncation level of the sum falls as the cutoff rises; a genuine singular value
does not move with the cutoff; and a residual limited by the working precision does not move with
the cutoff either, so the cutoff alone cannot tell the last two apart. Two cutoffs can share one
precision floor. A verdict is therefore taken at two or more cutoffs \emph{and} at two
precisions, and it is the precision that does the work.

An earlier draft of this section claimed that nothing had been seen between $10^{-2}$ and
$10^{-12}$ on this project's curves, and used that as a second reason to trust a reading. The
claim is false and the record contains the counterexample. At $X_0^{4130}(1)$ the smallest
singular value sits at $1.919\cdot10^{-3}$ at cutoffs $40$, $46$ and $50$, all at $50$ digits,
agreeing to eleven digits across them, and the count there is zero: not hyperelliptic. At $80$
digits, at the same cutoffs, it is $2.002\cdot10^{-10}$ with a jump of $1.4\cdot10^{9}$ above it:
one relation, hyperelliptic. A residual can therefore sit in the middle of the spectrum, be
perfectly stable in the cutoff, and be an artefact of the precision. $X_0^{6}(649)$ is the same
trap in the other direction: at $50$ digits the residual \emph{rose} from $7.7\cdot10^{-14}$ to
$3.8\cdot10^{-8}$ as the cutoff went from $40$ to $46$, which reads as a refusal, while at $90$
digits it is $1.5\cdot10^{-13}$ and $2.9\cdot10^{-16}$.

So the band is a description of the well-resolved cases, not a test, and the asymmetry matters:
the reading a precision floor fakes is the \emph{negative} one, since a floor puts a nonzero
number where a vanishing singular value belongs. A hyperelliptic reading is protected by the gap
it must show; a non-hyperelliptic reading is not protected by anything except a second precision,
and the code refuses to return one without it.

The test in both directions is the working precision. A genuine singular value does not notice it
(on the non-hyperelliptic control the smallest singular value agrees to twenty-nine digits between
thirty and fifty digits of precision), while a value limited by the precision moves with it: by a
factor of about a hundred at the first base where a floor was met, and by seven orders of
magnitude at $X_0^{4130}(1)$ above. So a residual is settled by raising the precision, never by
assumption, and the seventeen readings this project has taken of the open quotients are all at the
highest precision tried for each base.

The verdict is numerical. The project treats it as a strong indication that directs the exact
methods (a model, or an exact argument that excludes a hyperelliptic involution), not as a proof.

\section{The model}

Let $X=X_0^D(N)/W$ have genus three and suppose that an Atkin--Lehner element $u\notin W$ has
$X/\langle W,u\rangle$ of genus one. The canonical image of a hyperelliptic curve carries no
moduli, so the values of the differentials alone cannot give the equation; one derivative does.
With $x=u_2/u_1$ a coordinate on the canonical conic, the function $y=(dx/dz)/u_1(z)$ is odd
under the hyperelliptic involution and $y^2=f(x)$ with $f$ of degree eight. The derivative is
taken by Cauchy's formula on a circle of radius $\rho$ around each sample point, where the
trapezoidal rule converges like $(\rho/R)^n$, $R$ the distance to the nearest pole of $x$, and the
noise of the input is amplified only by $1/\rho$; the same circle gives the derivative from half
its points, which is the convergence check.

\subsection*{The derivative, from inside the sum}
A circle costs a dozen evaluations for one derivative, and it is not needed. The point $z$ enters
\eqref{eq:nelson} only through $\sigma(z)$, so the sum may be differentiated under the summation
sign, and the enumeration of short vectors --- which is the entire cost --- is then shared between
a value and its derivative. That the individual terms are not holomorphic in $z$ is no obstacle:
it is the $\partial_{\bar z}$ parts that cancel in the sum, not the $\partial_z$ ones, so
$\partial_z$ may be applied term by term. The matrices entering $X$ are real, so $X$ and $\bar X$
are the two linear functionals $\tfrac12(M_{11}+M_{22})\pm\tfrac{i}{2}(M_{12}-M_{21})$ of $M$, and
the same two applied to $\partial_z M=\sigma(z_1)^{-1}\varrho(b)\,\partial_z\sigma(z)$ give
$\partial_zX$ and $\partial_z\bar X$, whence
$\partial_z|X|=(\bar X\partial_zX+X\partial_z\bar X)/2|X|$. Each term also needs $W_k'$, which
comes from differentiating the Chebyshev series already fitted for $W_k$.

Against the circle at $X_0^{35}(1)/w_5$ the two derivatives agree to $10^{-14}$, which is the
circle's own accuracy, and the evaluation is $7.6$ times faster; over a whole model build the
saving is a factor $5.3$. Carrying no radius, the derivative inside the sum is also the more
accurate of the two, and the circle is kept as its check. What the values return is
$g_i=y_1y_2\overline{f_i(z_1)}f_i(z_2)$, which the factor $y_2$ makes non-holomorphic; the
holomorphic $h_i=g_i/y_2$ has
\[
  h_i'=\frac{\partial_zg_i+\tfrac{i}{2}h_i}{y_2},
\]
since $\partial_z\operatorname{Im}z=-i/2$.

\subsection*{The coordinate}
The hyperelliptic involution is central in the automorphism group, so $u$ descends to an
involution of the conic and the coordinate can be chosen with $u\colon x\mapsto -x$. Concretely,
$u$ fixes one of the three differentials and negates the other two; the conic has no cross term
between the fixed one and the other two (this is checked, and it would fail if the eigenvalue
convention were wrong), so it reads $q_{pp}e_p^2+(\text{binary form in }e_a,e_b)$, and factoring
the binary form gives $w_2^2=w_1w_3$ with $w_2=e_p$. Then $x=w_2/w_1$, the octic is even,
$f(x)=g(x^2)$, and the quotient $X/u$ is $Y^2=g(t)$. The evenness is read off the fit, not
assumed: the quartic $g$ in $t=x^2$ has five coefficients, is fitted on five sample points, and
must predict $y^2$ at the others.

\subsection*{The invariants}
The coordinate-free invariants of a binary octic have heights far beyond what the values carry.
The quartic $g$ with its two marked points $t=0,\infty$ has instead the three numbers
\[
A=\frac{c_0c_2}{c_1^2},\qquad B=\frac{c_1c_3}{c_2^2},\qquad C=\frac{c_2c_4}{c_3^2},
\]
unchanged by $t\mapsto\lambda t$ and by rescaling $g$, which fifteen digits settle at the
control $X_0^{35}(1)$ ($A=1939/952200$, $B=34500/76729$, $C=1939/5000$). Each is recognised by
the classical certainty criterion: a fraction $p/q$ within $\varepsilon$ of the computed real part
is accepted when $q^2\varepsilon<\tfrac12$, so that no other fraction of denominator at most $q$
lies as close, with $\varepsilon$ taken from the imaginary part of the computed number, which is
pure noise. Swapping the two roots of the binary form is $x\mapsto 1/x$, which reverses $g$ and
exchanges $A$ and $C$; both orderings are legitimate.

\subsection*{The twist}
$A,B,C$ fix $g$ up to $t\mapsto\lambda t$ and rescaling, so the curve is
\[
y^2=\kappa\,g_0(\mu x^2),\qquad g_0(t)=1+t+At^2+A^2Bt^3+A^3B^2Ct^4,
\]
with $\kappa,\mu\in\Q^\times$ mattering only modulo squares, and a complex-analytic computation
cannot see either: its coordinate differs from a rational one by a complex scalar, which moves
$\kappa$ and $\mu$ by complex squares. Both are solved rather than searched. The genus-one
quotient is $Y^2=\kappa g_0(t)$, $\mu$ having disappeared under $t\mapsto t/\mu$, so at each odd
prime $p$ of good reduction the point count of $Y^2=g_0(t)$ over $\F_p$ against the trace formula
for $X/\langle W,u\rangle$ gives the Legendre symbol $(\kappa\mid p)$, one linear equation over
$\F_2$ in the exponents of $\kappa$ over the primes of $2DN$ and the bad primes of $g_0$; the
system is heavily overdetermined and must have a unique solution. Then $\mu$ is read off the
curve itself the same way, $\kappa$ being known. The resulting curve is put in reduced minimal
form and its point counts over $\F_p$ and $\F_{p^2}$ compared with the trace formula.

\subsection*{When the coordinate is not rational}
The two fixed points of $u$ on the conic are the points $x=0,\infty$; they are either two rational
points or a conjugate pair over a quadratic field $K=\Q(\sqrt d)$. In the second case the
coordinate with $u\colon x\mapsto -x$ exists only over $K$, the quartic $g$ has coefficients in
$K$, and the nontrivial element $\sigma$ of $\mathrm{Gal}(K/\Q)$ acts as $x\mapsto a/x$ for some
$a\in\Q^\times$, so that $B$ is rational while $A$ and $C$ are conjugate. This is what the
values show at $X_0^{6}(323)/\langle w_6,w_{19},w_{34}\rangle$: $B$ recognised, $A$ and $C$
complex conjugates to sixteen digits, with the conic and the predictions as clean as anywhere.

The descent goes as follows. $K$ is read off the arithmetic: the fixed points of $u$ on $X$ are
CM points, their fields of definition contain $K$, and $K$ is the quadratic subfield they have in
common (at the base above, one discriminant, $-7752$, carries all four points, over a quartic
field whose quadratic subfield is $\Q(\sqrt{-2})$). $A$ is then recognised in $K$ by simultaneous
rational approximation of its two coordinates, which certifies a denominator up to about
$(2\varepsilon^2)^{-1/3}$ at absolute accuracy $\varepsilon$. At the base above the data give
$\varepsilon=4.4\cdot10^{-15}$, hence denominators up to $2.97\cdot10^9$, and the $A$ found has
denominator $1\,262\,027\,560$, inside that bound; a fraction of larger height would be missed,
not misread. Write $F(x)=g_0(\mu x^2)$ with
$g_0(t)=1+t+At^2+A^2Bt^3+A^3B^2Ct^4$ and $\mu\in K^\times$. The identity
$t^4g_0(1/t)=A^3B^2C\,g_0^\sigma(t/ABC)$ turns the condition that the branch divisor be
Galois-stable, $F^\sigma(x)=c\,x^8F(a/x)$, into the norm equation
$N_{K/\Q}(\mu)=1/(a^2B\,N_{K/\Q}(A))$, solvable exactly when $B\,N(A)$ is a norm from $K$
(Hilbert~90); when it is not, that is the result. The solution $\mu$ is defined up to norm-one
elements $\sigma(\gamma)/\gamma$, and $(\mu\sigma(\gamma)/\gamma,\,aN(\gamma))$ describes the
same curves as $(\mu,a)$: the forms over $\Q$ of the marked curve are the twists by the involution
$u$ over $\Q(\sqrt e)$, $e\in\Q^\times/\Q^{\times2}$, which appear as $(\mu/e,\,a=e)$, together
with the quadratic twists by $r$. The quotient by $u$ is the same for every $e$. The descent
datum of $y^2=\kappa F_e(x)$ is $(x,y)\mapsto(e/x,\,\epsilon y/x^4)$ with
$\epsilon^2=\sigma(\kappa)c/\kappa$ and $N(\epsilon)=e^4$; taking $\epsilon=e^2$, Hilbert~90 gives
$\kappa_0=1+c/e^4$, and every valid scalar is $\kappa_0$ times a rational number $r$ times a
square. Nothing here needs $e$ to be a norm from $K$. When $e$ is a norm, $c\in K$ with $N(c)=e$
splits the cocycle through $x=c(1+\sqrt d\,s)/(1-\sqrt d\,s)$, $s$ rational, and the curve is
$y^2=rP(s)$ over $\Q$ with $u$ acting as $s\mapsto 1/(ds)$; when $e$ is not a norm, the conic
$X/\iota$ is $X^2-dY^2=eZ^2$, which has no rational point, the curve has no model $y^2=f(s)$ over
$\Q$, and the $K$-model with its datum is the result.

Both $e$ and $r$ are read off point counts. At a prime $p$ split in $K$ the curve over $\F_p$ is
the $K$-model reduced at a prime above $p$, so $e$ is found first, as the class whose counts match
for one of the two quadratic twists, by a linear system over $\F_2$ in its exponents over the
primes of $2DNd$ and of $\mu$; then $r$ from the same primes on $y^2=\kappa_0F_e$. The factor $d$
in each is invisible at split primes, where $d$ is a square; at an inert prime the $\F_p$-points
of the curve are the points $(x,y)$ over $\F_{p^2}$ with $x^{p+1}=e$, $y^2=F(x)$ and
$y^{p-1}=\epsilon/x^4$, the factor $d$ in $r$ is the sign of $\epsilon$, and these counts settle
both. The procedure is checked on curves whose equations are known: a synthetic octic with the
symmetry $s\mapsto 1/(ds)$ over $\Q(\sqrt{-2})$ and $\Q(\sqrt2)$, moved to the $K$-coordinate and
recovered up to isomorphism over $\Q$; the exact model of $X_0^{1326}(1)/\langle w_{17},w_{78}\rangle$
with its two further involutions, whose fixed points are conjugate over $\Q(\sqrt{-2})$ and over
$\Q(\sqrt2)$, recovered up to isomorphism over $\Q$ with the twists read off the trace formula, the
quotients having Jacobians $1326c1$ and $1326b2$; and the model of
$X_0^{210}(13)/\langle w_2,w_{21},w_{65}\rangle$ with its involution over $\Q(i)$.

At $X_0^{6}(323)/\langle w_6,w_{19},w_{34}\rangle$ the outcome is the second kind: $e$ is not a
norm from $\Q(\sqrt{-2})$, the conic has no rational point, and the curve is given by its model
over $\Q(\sqrt{-2})$ with the datum $\sigma(x)=e/x$. The four checks pass: the fixed points of
$u$ on the model are defined over a quartic field isomorphic to the field of definition of the CM
points of discriminant $-7752$; the Jacobian of the descended quotient is the elliptic curve
$y^2+xy=x^3-139x-271$ of conductor $1938$, with the Fourier coefficients of the $W$-fixed newform
with $w_2=+1$; the $W$-fixed newforms split $1+2$; and the numerators of the zeta functions at the
split primes agree with the trace formula, as do the counts over $\F_p$ at the inert primes through
the datum. The same label applies as to every model here: numerical, confirmed by point counts.

\subsection*{Where the fit is conditioned, and where it is not}
The quartic is determined at its large-$t$ end by the sample points with large $\lvert x\rvert$
and at its small-$t$ end by the others. When the moduli $\lvert x\rvert$ at the sample points span
less than a factor of about five, the fit is ill conditioned and the invariants come out to far
fewer digits than the values carry, with the predictions at $10^{-6}$ while the conic is at
$10^{-15}$; this happened at several bases. The remedy is to add sample points where
$\lvert x\rvert$ is large, which a coarse evaluation at a small cutoff finds, since $x$ is read off
the conic.

\section{The checks, and their status}

Point counts determine the isogeny class of the Jacobian, not the curve, so four further checks
are made, each of which the isogeny class cannot supply.

\begin{enumerate}
\item[(a)] The four fixed points of $u$ on $X$ are the images of the fixed points of the elements
of the coset $uW$, CM points whose discriminants and multiplicities are known; their total over
$\#W$ must be $4$, as Riemann--Hurwitz demands for a double cover of a genus-one curve by a
genus-three curve. On the model they are the points over $x=0$, defined over $\Q(\sqrt\kappa)$,
and over $x=\infty$, defined over $\Q(\sqrt{\text{leading coefficient}})$; the fields of definition
of the CM points on $X$ are computed independently, by Shimura reciprocity, and must match pair by
pair. A discriminant carrying four points of $X$ defined over one quadratic field contributes that
field twice.
\item[(b)] The Jacobian of $X/\langle W,u\rangle$ has conductor $DN$ and the Fourier coefficients
of the $W$-fixed newform of level $DN$ with quaternionic $u$-eigenvalue $+1$, at every good prime
below a hundred.
\item[(c)] The $W$-fixed newforms of level $DN$ split $1+2$ under $u$.
\item[(d)] The numerator of the zeta function of the model over $\F_p$, for good $p\le60$, equals
the one assembled from the trace formula's counts over $\F_p,\F_{p^2},\F_{p^3}$, which determine
it; likewise for the quotient.
\end{enumerate}

\subsection*{Status}
The construction is numerical and the checks are arithmetic consequences of the isogeny class
and of the Galois structure of the branch locus; none of them proves the model. The strongest
evidence available is agreement with a model obtained exactly. At $X_0^{1326}(1)/\langle
w_{17},w_{78}\rangle$, whose Borcherds level lies beyond the reach of the project's lattice-point
enumerator, the double cover of the elliptic curve $X/\langle W,w_2\rangle$ branched at the CM
points prescribed by the heights of Heegner points gives the exact model
\[
y^2=272x^8-72x^6-55x^4-18x^2+17,
\]
and the model recognised from the values of the three differentials at fourteen sample points,
with $u=w_2$, is isomorphic to it over $\Q$.

Two further comparisons with that construction, which has since reached more of these quotients,
say what the numerical route is worth where it does not finish. At
$X_0^{6}(185)/\langle w_2,w_{15},w_{111}\rangle$ the fit is clean but the certainty criterion
refuses two of the three invariants; against the exact model
$y^2=2025x^8+4716x^6-122x^4-5620x^2-999$ the three computed invariants agree to relative
$2\cdot10^{-12}$, $4\cdot10^{-12}$ and $2\cdot10^{-12}$, with the measured errors matching the
noise estimates the criterion reads off the imaginary parts. The refusals are correct rather than
conservative: $q^2\varepsilon$ is about $15$ for the invariant of denominator $15\,792\,200$ and
$0.8$ for the one of denominator $3721$, both refused, against $0.003$ for the one of denominator
$137\,288$, which was accepted. So at that base the limit is the height of the invariants against
the conditioning of the fit, and nothing else; the margin at $3721$ is thin, which measures how
little more accuracy would settle it. Which of the two outer invariants is which is not
determined by the fit -- exchanging the roots of the binary form is $x\mapsto1/x$, which reverses
the quartic -- so a comparison must allow either ordering, and the two runs made of this base
chose differently. And the pointless conic of the previous
section is not peculiar to one base: the exact route finds the same phenomenon at
$X_0^{6}(187)$ and $X_0^{22}(51)$, hyperelliptic quotients with no model $y^2=f(x)$ over $\Q$.

\section{A genus-two quotient, and a model over $\Q$}\label{sec:genus2}

Let $X=X_0^D(N)/W$ have genus two. The canonical map is then the hyperelliptic map itself, so no
marked involution is needed and the equation comes from the differentials alone: with
$\omega_i=h_i\,dz$ a basis of them,
\[
  x=\frac{h_2}{h_1},\qquad y=\frac{dx/dz}{h_1},
\]
since $dx/y$ and $x\,dx/y$ are a basis of the differentials of $y^2=F(x)$. The relation
$y^2=F(x)$ is therefore \emph{linear} in the seven coefficients of $F$: seven sample points
determine it, and the others test it. Equivalently, with $\mathcal W=h_1h_2'-h_1'h_2$ the
Wronskian --- a form of weight six vanishing exactly at the six Weierstrass points --- it reads
$\mathcal W^2=\sum_k c_kh_2^kh_1^{6-k}$, a relation among forms of weight twelve.

\subsection*{A coordinate in which the equation is rational}
$F$ comes out with complex coefficients, $h_1,h_2$ being some complex basis: the curve is right
and its equation is not over $\Q$. Recognising the absolute invariants instead does not work,
their heights lying far beyond what the values carry --- at the control they have numerators of
order $10^{17}$ and would need some twenty-six digits. So the coordinate is fixed instead. A CM
point defined over $\Q$ is a rational point of $X$, so in a coordinate carrying three of them to
$0,1,\infty$ every rational point has rational coordinate, and the sextic then has rational
coefficient ratios --- of small height, being an honest equation of the curve. The rational CM
points that fixed nothing must come out rational in that coordinate too, which is a second test
alongside the sample points held back.

\subsection*{A CM point at which the coordinate is $0/0$}
The $h_i$ are pullbacks of the differentials of the quotient, and a pullback vanishes wherever the
covering map ramifies. So at a point fixed by some $w_m\in W$ --- a CM point of discriminant
$-4m$, or $-m$ when $m\equiv3\pmod 4$ --- \emph{both} $h_1$ and $h_2$ vanish, $x=h_2/h_1$ is
$0/0$, and what is computed there is a quotient of two rounding errors. At a full star every $w_m$
lies in $W$, so the case cannot be avoided. An Atkin--Lehner involution has order two, so both
forms vanish to order exactly one and the limit is $h_2'/h_1'$, which the derivative of the
previous section supplies at no extra cost.

At $X_0^{2002}(1)^*$ this is visible without any model. Seven of the nine rational CM points share
one line through the origin to $6\cdot10^{-16}$ --- the real locus is a line in this normalisation,
and a rational point must lie on it --- while the two that miss it, by $1.6\cdot10^{-2}$ and
$0.85$, are exactly those whose discriminants $-11$ and $-91$ are fixed by $w_{11}$ and $w_{91}$.
There $|h_1|\approx10^{-16}$, against $10^{-1}$ at an ordinary point of the same run.

\subsection*{The twist, and the control}
Mestre's construction is not needed, the sextic being obtained directly; what the normalisation
discards is a single scalar, so the curve is determined up to a quadratic twist. The twist is read
off the counts: the right curve has $\#C(\F_p)=p+1-\operatorname{Tr}a_p(f)$, so at each good prime
the untwisted model gives one of $p+1\mp\operatorname{Tr}a_p$, the pattern of signs is the
quadratic character of the twisting discriminant, and that discriminant can ramify only where the
model is bad --- finitely many characters, cut out by Legendre symbols over $\F_2$. The newform
orbit is identified by $|\operatorname{Tr}a_p|$ rather than by Atkin--Lehner signs, which is
convention-free and at the same time checks the isogeny class.

At $X_0^{35}(1)/w_5$, whose model is known independently from Guo and Yang, the route returns
\[
  y^2=-2401u^5+4900u^4-3878u^3+1380u^2-u
\]
in the coordinate sending the CM points of discriminant $-28,-35,-7$ to $0,1,\infty$, with
twisting constant $-1$. The reduced minimal model then agrees with the stored one equation for
equation, the absolute invariants agree \emph{as rationals} rather than numerically, and the counts
agree with the newform at all forty-three good primes below two hundred. The route's own two tests
pass as well: the three sample points held back are predicted to $2\cdot10^{-17}$, and the rational
CM points that fixed nothing come out rational.

This is worth separating from the method of the previous section. There the curve is pinned only up
to its isogeny class until something external selects it, which is the difficulty a surface with
several principal polarizations raises: all of them share an unpolarized Jacobian, hence an
$L$-function, hence every point count. Here the curve is read off its own canonical differentials,
no choice is made, and the counts are used only for the twist --- which they do determine.

\section{Three traps}

\subsection*{Precision must rise with the cutoff}
Raising the cutoff stops helping once the truncation error falls below the rounding error of the
sum. At $X_0^{1326}(1)$ with thirty digits the residual of Noether's count was the same at cutoffs
$46$ and $50$ while the rest of the spectrum agreed to twelve digits, which is exactly what a
genuine small singular value would look like; raising the precision to fifty digits at cutoff
$46$ lowered it by a factor of a hundred, and the curve is hyperelliptic, as the exact model
confirms. The sum there has some sixty million terms; above roughly ten million terms the
working precision must rise with the cutoff. Two other explanations were tested first and ruled
out: the fixed scaling by which the Gram matrix is rounded before the enumeration changes no digit
(the scaling only decides which vectors are offered, and each is then accepted or rejected at full
precision), and widening the margin on the Fourier coefficients changes no digit either.

\subsection*{The Eichler order is not deterministic across processes}
Magma's construction of an Eichler order of level $N$ returns a different (conjugate) order from
one process to the next, with the algebra's standard form and the real splitting unchanged. Values
computed against different orders are values at different points of the curve, with different
normalising constants, so a computation split over processes gives chunks that each lie on a
canonical conic and share none. The order is therefore fixed once, through its $\Z$-basis in the
standard basis of the algebra, and rebuilt in every process; with the same order the split
computation reproduces the single-process invariants of the control to every digit. The same
care is needed for the Galois conjugates of a newform orbit: the generator of the coefficient
field, and with it the order in which the conjugates are listed, is not the same from one process
to the next either, so the conjugates are ordered by their coefficients.

\subsection*{A denominator bound must come from the accuracy, not the precision}
\texttt{BestApproximation}$(\alpha,Q)$ returns a fraction within $1/Q^{2}$ of \emph{any} real
whatever, so a bound taken from the working precision rather than from the accuracy achieved
guarantees a spurious fit. Fitted at twelve digits, one coefficient of the control's sextic came
back as $-6227397235/7868552463$ in place of $-277/350$, with a smaller residual than the truth.
Both the tolerance and the bound are therefore read off the accuracy that the fit itself measures
at the points held back: $Q$ is chosen so that $1/Q^{2}$ is four orders of magnitude looser than
that accuracy, and the factor by which each coefficient beats $1/Q^{2}$ is reported next to it. At
the control those factors run from $10^{5}$ to $10^{14}$. The same caution applies to a vanishing
coefficient, which shows up at the size of the accuracy and not below it: judged against the
working precision it looks like a failure, and the run then reports no rational sextic on data that
is in fact exact.

\section{What is implemented}

The package \texttt{NelsonShimizu.m} provides: the Eichler order serialised and rebuilt; the
basis of $W$-fixed differentials, newforms of level exactly $DN$ together with the oldform
eigenvectors $v_\chi$, with the dimension check; the evaluation of the basis at a pair of points
by \eqref{eq:nelson} and \eqref{eq:pullback}, with the cutoff and precision as parameters; the
direct measurement of the $W$-eigenvalues of every column; Noether's count with the reading rule
above; the values on circles, their storage, and the model step through the elliptic quotient;
the values together with their derivatives in $z$, and from them the hyperelliptic coordinate and
$dx/dz$ that a genus-two equation needs; and the four checks. Every docstring says that its result
is numerical. The two drivers \texttt{tools/genus2model.m} and \texttt{tools/genus2twist.m} carry
out the genus-two route of \S\ref{sec:genus2} and the twist that completes it.

\begin{thebibliography}{9}
\bibitem{JL} H.~Jacquet and R.~P. Langlands, \emph{Automorphic forms on $GL(2)$}, Lecture Notes
in Mathematics \textbf{114}, Springer (1970).
\bibitem{Nelson} P.~D. Nelson, \emph{Evaluating modular forms on Shimura curves}, Math. Comp.
\textbf{84} (2015), 2471--2503.
\bibitem{Voight} J.~Voight, \emph{Quaternion algebras}, Graduate Texts in Mathematics
\textbf{288}, Springer (2021).
\bibitem{VW} J.~Voight and J.~Willis, \emph{Computing power series expansions of modular forms},
in: Computations with modular forms, Contrib. Math. Comput. Sci. 6, Springer (2014), 331--361.
\end{thebibliography}

\end{document}
